\documentclass[10pt,a4paper]{amsart}
\usepackage[margin=19mm]{geometry}
\usepackage{amsmath,amssymb,mathtools}
\usepackage[scr=rsfs,cal=euler]{mathalfa}
\usepackage{xfrac}
\usepackage[unicode,hidelinks,bookmarks=false]{hyperref}
\newcommand{\ie}{i.e.\ }
\newcommand{\cf}{cf.\ }
\newcommand{\resp}{respectively\ }
\newcommand{\Subsection}{\subsection}
\providecommand{\nobreakdash}{\nobreak\hbox{-}}
\setlength{\emergencystretch}{2em}
\multlinegap0pt
\usepackage[shortlabels]{enumitem}
\usepackage{amssymb}
\usepackage{tikz}
\usepackage{dsfont}
\usepackage{float}
\newtheorem{lem}{Lemma}
\newcommand{\De}{\Delta}
\newcommand{\ba}{\begin{equation} \begin{aligned}}
\newcommand{\bs}[1]{\boldsymbol{#1}}
\newcommand{\ea}{\end{aligned} \end{equation}}
\newcommand{\rhob}{{\bar{\rho}}}
\title{Maximally heavy dynamics in the causal diamond}
\author{David Poland, Gordon Rogelberg}
\date{Source: arXiv:2603.28853v1 (CC BY 4.0)}
\begin{document}
\maketitle

\noindent\emph{Source and licence.} Maximally heavy dynamics in the causal diamond. Version arXiv:2603.28853v1 (CC BY 4.0), distributed by arXiv under the Creative Commons Attribution 4.0 International licence, http://creativecommons.org/licenses/by/4.0/, as recorded in the arXiv licence metadata for this exact version and shown on the abstract page https://arxiv.org/abs/2603.28853v1. Full licence notice: \texttt{licenses/arxiv-2603.28853-CC-BY-4.0.txt}.\par\smallskip

\noindent\textit{Editorial scope.} Counted unit: the article's lem sharpmomentbound (source0.tex lines 1895--1901). The article's complete original proof of it is delivered with it (source0.tex lines 1902--1923, one \texttt{proof} environment of 22 lines). No mathematics is written, repaired or re-derived: the statement and the proof are the article's own lines, in the article's own notation, and no retained line is substituted or re-flowed.  No further source-local context is needed: every cross-reference of every retained line resolves inside the two delivered spans.  Nothing was omitted from any retained span, so the package carries no omission record.  No retained line contains a citation, so the wrapper carries no bibliography and no bibliography entry.  No retained line contains a figure environment, an image-inclusion command or drawing code, so no image is delivered and the package declares no asset dependency.  Editorial formatting only, in addition: an AMS \texttt{amsart} preamble that restates the article's own declarations and macros used by the retained lines, namely the environments \texttt{lem} on separate counters, together with the standard packages those lines need.  The retained environments print in this document as Lemma 1 in order of appearance, whereas the article prints the same items with the numbering of its own full document.  The complete native source of the article as distributed in the arXiv e-print for this exact version is bundled unchanged as \texttt{sources/197472/source0.tex}.
\begin{lem}[Sharp bound on the first moment in the causal diamond]
\label{sharpmomentbound}
The first normalized moments satisfy the sharp bounds
\ba
\omega_{1,0}(\bs{\rho}) \leq \De \frac{1+\rho}{1-\rho} \quad \text{and} \quad \omega_{0,1}(\bs{\rho}) \leq \De \frac{1+\rhob}{1-\rhob} \quad \forall \bs{\rho} \in (0,1)^2.
\ea
\end{lem}

\begin{proof}
The moments satisfy $\omega_{j,k}(\bs{\rho}) = \omega_{k,j}(\bar{\bs{\rho}})$, so it suffices to check for $(j,k) = (1,0)$ as long as the bound holds for all $\bs{\rho} \in (0,1)^2$. We write $\bs{\rho} = (e^{-\bs{s}}) = (e^{-s},e^{-\bar{s}})$ and take the logarithm of the crossing equation
\ba
\log(\mathcal{G}(\bs{s})) = \De \log(|\bs{\chi}(\bs{s})|) + \log( \mathcal{G}( \hat{\bs{s}} )),
\ea
where 
\ba
\hat{\bs{s}} = \left(- \log(\hat{\bs{\rho}}(\bs{s}))\right) = \left(- 2 \log\left(\frac{1-\sqrt{e^{-s}}}{1+\sqrt{e^{-s}}}\right) ,- 2 \log\left(\frac{1-\sqrt{e^{-\bar{s}}}}{1+\sqrt{e^{-\bar{s}}}}\right) \right)
\ea
and 
\ba
|\bs{\chi}(\bs{s})| = \left(  \frac{16 e^{-(s+\bar{s})}}{(1-e^{-s})^2 (1-e^{-\bar{s}})^2}\right).
\ea
We now take a derivative of both sides with respect to $-s$ and evaluate at $\bs{s} = \left(-\log(\bs{\rho}) \right)$ to compute the first moment $\omega_{1,0}(\bs{\rho})$ on the LHS
\ba
\omega_{1,0}(\bs{\rho}) = \De \partial_{-s} \log(|\bs{\chi}(\bs{s})|)|_{\bs{s} = -\log(\bs{\rho})} + \left(\partial_{-s} \hat{\rho}(s) \right)|_{s = -\log(\rho)} \omega_{1,0}(\hat{\bs{\rho}}).
\ea
By the positivity of the support of the OPE measure, $\omega_{1,0}(\bs{\rho})$ is positive for all points in the causal diamond, as is $\omega_{1,0}(\hat{\bs{\rho}})$. It is also easy to check that $\partial_{-s} \hat{\rho}(s) = \frac{2 \sqrt{e^{-s}}}{e^{-s} -1} \leq 0$ for all points in the causal diamond, so the second term on the RHS is always negative. Thus, we have
\ba
\omega_{1,0}(\bs{\rho}) &= \De \partial_{-s} \log(|\bs{\chi}(\bs{s})|)|_{\bs{s} = -\log(\bs{\rho})} + \left(\partial_{-s} \hat{\rho}(s) \right)|_{s = -\log(\rho)} \omega_{1,0}(\hat{\bs{\rho}})\\& \leq \De \partial_{-s} \log(|\bs{\chi}(\bs{s})|)|_{\bs{s} = -\log(\bs{\rho})}  = \De \frac{1+\rho}{1-\rho}.
\ea
\end{proof}

\end{document}
